% context: markscheme (official solutions) for 1-7Q_Measurement
% topic: significant figures, rounding, scientific notation, reading a ruler,
%        unit conversion, average speed
% convention: "=" joins exact equalities; "\approx" introduces a rounded value;
%        long decimals use \num{} so the digits are grouped in threes
\documentclass[12pt, twoside]{article}
%\documentclass[12pt, twoside]{article}
\usepackage[letterpaper, margin=1in, headsep=0.2in]{geometry}
\setlength{\headheight}{0.6in}
%\usepackage[english]{babel}
\usepackage[utf8]{inputenc}
\usepackage{microtype}
\usepackage{amsmath}
\usepackage{amssymb}
%\usepackage{amsfonts}
\usepackage[nomessages]{fp} %\FPeval{\var-name}{2*sin(pi/6)}
\usepackage{siunitx} %units in math. eg 20\milli\meter
\usepackage{yhmath} % for arcs, overparenth command
\usepackage{tikz} %graphics
\usetikzlibrary{quotes, angles, arrows, arrows.meta}
\usepackage{graphicx} %consider setting \graphicspath{{images/}}
\usepackage{parskip} %no paragraph indent
\usepackage{enumitem}
\usepackage{multicol}
\usepackage{venndiagram}

\usepackage{fancyhdr}
\pagestyle{fancy}
\fancyhf{}
\renewcommand{\headrulewidth}{0pt} % disable the underline of the header
\raggedbottom
\hfuzz=2mm %suppresses overfull box warnings

\usepackage{hyperref}
\usepackage{float}
\setlength{\parskip}{0.3\baselineskip}
\sisetup{reset-text-series=false, reset-math-version=false, text-series-to-math=true} % \num inherits bold

\title{Physics}
\author{Chris Huson}
\date{October 2026}

\fancyhead[LE]{\thepage}
\fancyhead[RO]{\thepage}
\fancyhead[LO]{La Scuola d'Italia / Huson / Physics \\* 2 October 2026}

\begin{document}

\subsubsection*{1.7 Quiz Solutions: Measurement}

\emph{Notation: \,$=$\, joins values that are exactly equal; \,$\approx$\, introduces a
rounded value. Each conversion runs as one chain of unit factors and is evaluated once.}

\begin{enumerate}[itemsep=0.25cm]

\item Significant digits.
  \begin{multicols}{3}
    \begin{enumerate}[itemsep=0.1cm]
      \item \textbf{3}
      \item \textbf{3}
      \item \textbf{2}
      \item \textbf{4}
      \item \textbf{1}
      \item \textbf{3}
    \end{enumerate}
  \end{multicols}
  \emph{(a) Leading zeros do not count; the trailing zero after the decimal point does.
  (b) A zero between nonzero digits always counts. (e) \num{0.00009} has a single
  measured digit, the 9.}

\item Round to three significant figures.
  \begin{enumerate}[itemsep=0.15cm]
    \item $\sqrt{13} = 3.6055512\ldots \approx \mathbf{3.61}$
    \item $8848.86~\text{m} \approx \mathbf{8850~\textbf{m}}$, better written
      $\mathbf{8.85 \times 10^{3}~\textbf{m}}$

      \emph{The third figure is the 4 in the tens place; the 8 after it rounds it up.}
  \end{enumerate}

\item Scientific notation.
  \begin{enumerate}[itemsep=0.15cm]
    \item $9{,}460{,}730{,}000{,}000~\text{km} = 9.46073 \times 10^{12}~\text{km}
      \approx \mathbf{9.46 \times 10^{12}~\textbf{km}}$
    \item \num{0.0000000000000000001602}~C $= 1.602 \times 10^{-19}$~C
      $\approx \mathbf{1.60 \times 10^{-19}~\textbf{C}}$

      \emph{Six groups of three zeros is 18 zeros, so the 1 is in the 19th decimal
      place: the exponent is $-19$. Write 1.60, not 1.6 --- the trailing zero is the
      third significant figure.}
  \end{enumerate}

\item Ordinary decimal numbers.
  \begin{enumerate}[itemsep=0.15cm]
    \item $7.0 \times 10^{-5}~\text{m} = $ \textbf{\num{0.000070}~m}

      \emph{Keep the trailing zero: $7.0$ has two significant figures and so must the
      decimal form. 0.00007 m loses one.}
    \item $5.46 \times 10^{7}~\text{km} = \mathbf{54{,}600{,}000~\textbf{km}}$
  \end{enumerate}

\item Reading the ruler. The bar ends between the 8.3 and 8.4 cm marks, a little past
  halfway.
  \begin{enumerate}[itemsep=0.15cm]
    \item $\mathbf{8.36~\textbf{cm}}$ --- the 8.3 is certain and the final 6 is the
      estimated digit. \emph{Accept 8.35 to 8.37.}
    \item \textbf{Three significant figures.}
    \item $8.36~\text{cm} \times \dfrac{10~\text{mm}}{1~\text{cm}}
      = \mathbf{83.6~\textbf{mm}}$
  \end{enumerate}

\item Convert.
  \begin{enumerate}[itemsep=0.15cm]
    \item $508~\text{m} \times \dfrac{3.28~\text{ft}}{1~\text{m}} = 1666.24~\text{ft}
      \approx \mathbf{1670~\textbf{ft}}$, or $\mathbf{1.67 \times 10^{3}~\textbf{ft}}$

      \emph{The published height is 1667 ft, so this checks.}
    \item $310~\dfrac{\text{km}}{\text{h}} \times \dfrac{1000~\text{m}}{1~\text{km}}
      \times \dfrac{1~\text{h}}{3600~\text{s}} = 86.111\ldots
      \approx \mathbf{86.1~\textbf{m/s}}$
  \end{enumerate}

\item Cheetah.

  $29~\dfrac{\text{m}}{\text{s}} \times \dfrac{3.28~\text{ft}}{1~\text{m}}
  \times \dfrac{1~\text{mi}}{5280~\text{ft}} \times \dfrac{3600~\text{s}}{1~\text{h}}
  = 64.854\ldots \approx \mathbf{65~\textbf{mi/h}}$

  \emph{Two significant figures, set by the 29 m/s. Accept 64.9 mi/h for full credit on
  the conversion; the sig-fig point is worth discussing when the quiz is returned.}

\end{enumerate}

\subsubsection*{Early finisher}

\emph{Two steps: (1) the rate relationship, rate $\times$ time $=$ quantity or a
rearrangement; (2) convert that result to the units asked for, carrying the unrounded
value. Convert first when it simplifies Step 1; metric changes are done in the head.}

\begin{enumerate}[resume, itemsep=0.1cm]

\item Radio command to Mars. \quad $t = \frac{d}{r}$; \, $d$ is in km, so first
  \mbox{$3.00 \times 10^{8}$ m/s $= 3.00 \times 10^{5}$ km/s.}

  Step 1: \quad $t = \dfrac{5.46 \times 10^{7}~\text{km}}{3.00 \times 10^{5}~\text{km/s}}
  = 182~\text{s}$

  Step 2: \quad $182~\text{s} \times \dfrac{1~\text{min}}{60~\text{s}} = 3.0333\ldots
  \approx \mathbf{3.03~\textbf{min}}$

  \emph{That is at closest approach, and up to about 24 minutes when Mars is far ---
  which is why a Mars rover cannot be driven with a joystick.}

\item Widening of the Atlantic. \quad $t = 2026 - 1492 = 534$ yr.

  Step 1: \quad $d = r \cdot t = 2.5~\dfrac{\text{cm}}{\text{yr}} \times 534~\text{yr}
  = 1335~\text{cm}$

  Step 2: \quad $1335~\text{cm} = 13.35~\text{m} \approx \mathbf{13~\textbf{m}}$

  \emph{Two significant figures, set by ``about 2.5 cm''; the 534 years is a count and
  does not limit the answer. About the length of a city bus --- in five centuries.}

\item Shower.

  Step 1: \quad quantity $= r \cdot t = 2.5~\dfrac{\text{gal}}{\text{min}} \times
  8.0~\text{min} = 20~\text{gal}$

  Step 2: \quad $20~\text{gal} \times \dfrac{3.785~\text{L}}{1~\text{gal}} = 75.7~\text{L}
  \approx \mathbf{76~\textbf{L}}$

  \emph{Two significant figures: both 2.5 gal/min and 8.0 min have two.}

\item Water per person per day.

  Step 1: \quad $r = \dfrac{1.0 \times 10^{9}~\text{gal/day}}{8.5 \times 10^{6}~\text{people}}
  = 117.647\ldots$ gal per person per day

  Step 2: \quad $117.647\ldots~\text{gal} \times \dfrac{3.785~\text{L}}{1~\text{gal}}
  = 445.29\ldots~\text{L} \approx \mathbf{4.5 \times 10^{2}~\textbf{L}}$
  (450 L)

  \emph{Both data values have two significant figures, so the answer gets two; 445 L
  overstates the precision. Full credit for the calculation, the sig-fig mark only for two.}

\end{enumerate}
\end{document}
